% preamble.tex
% -----------------------------------------------------------------------
% "Aprendo a leer" -- cuaderno diario de lectura (español, A4).
% Todo el aparato tipográfico vive aquí: paquetes, colores, estilos de
% caja y las macros que usa cada página diaria. Requiere LuaLaTeX
% (fontspec) -- no pdflatex puro: la letra es Andika (SIL, OFL,
% fonts/andika/), pensada para primeros lectores -- "a"/"g" de un solo
% piso y "b"/"d"/"p"/"q" bien diferenciadas, a diferencia de Latin
% Modern Sans (ver notes/02-revision-y-plan.md, punto 5). fontspec ya
% cubre UTF-8 y la codificación de fuente por su cuenta -- no hacen
% falta fontenc ni inputenc, y de hecho estorban bajo LuaLaTeX.
% -----------------------------------------------------------------------

\usepackage{fontspec}
\setmainfont{Andika}[
  Path = fonts/andika/,
  Extension = .ttf,
  UprightFont = *-Regular,
  BoldFont = *-Bold,
  ItalicFont = *-Italic,
  BoldItalicFont = *-BoldItalic,
]
\setsansfont{Andika}[
  Path = fonts/andika/,
  Extension = .ttf,
  UprightFont = *-Regular,
  BoldFont = *-Bold,
  ItalicFont = *-Italic,
  BoldItalicFont = *-BoldItalic,
]

\usepackage[spanish]{babel}
% babel-spanish makes " an ACTIVE character (shorthand for <</>> and other
% combinations) -- a literal " typed in prose (dialogue, "quoted speech")
% collides with it and garbles silently, worst right before a capital
% letter ("E -> a stray accent, e.g. "Estoy" rendered as ".Ëtoy"). Plain
% \shorthandoff{"} here in the preamble does NOT work and fails silently
% -- babel re-activates its shorthands for the main language at
% \begin{document}, AFTER the preamble has already run, so it wins.
% \AtBeginDocument defers the command to run right after that, which is
% the only point it sticks -- verified against a minimal standalone file
% both ways before trusting either.
\AtBeginDocument{\shorthandoff{"}}

\usepackage[a4paper,top=18mm,bottom=16mm,left=22mm,right=18mm]{geometry}

\usepackage{xcolor}
\usepackage[most]{tcolorbox}
\usepackage{tabularx}
\usepackage{booktabs}
\usepackage{enumitem}
\usepackage{tikz}
\usepackage{multicol}

\pagestyle{empty}
\setlength{\parindent}{0pt}
\setlength{\parskip}{2pt}

% --- cadenas de texto (español; centralizadas, ver lang/es.tex) -------
% lang/es.tex
% -----------------------------------------------------------------------
% Todas las cadenas de texto que aparecen en el cuaderno, en un solo
% sitio. Cambiar una palabra aquí la cambia en las 260 páginas a la vez,
% sin tocar el contenido generado.
% -----------------------------------------------------------------------

\newcommand{\lblFecha}{Fecha}
\newcommand{\lblNotas}{Notas}
\newcommand{\lblDia}{Día}
\newcommand{\lblSemana}{Semana}
\newcommand{\lblTrimestre}{Trimestre}
\newcommand{\lblRespuesta}{Respuesta}

\newcommand{\lblInstruccionLectura}{Lee en voz alta, despacio, señalando cada palabra}

\newcommand{\lblDibuja}{Dibuja}
\newcommand{\lblCompleta}{Completa}
\newcommand{\lblCopia}{Copia}
\newcommand{\lblResponde}{Responde}
\newcommand{\lblRelaciona}{Relaciona}
\newcommand{\lblAdivina}{Adivina}
\newcommand{\lblCrea}{Crea}
\newcommand{\lblRepasa}{Repasa}

\newcommand{\lblInstruccionCopia}{Copia la frase de hoy, tres veces, una en cada línea.}
\newcommand{\lblInstruccionCompleta}{Convierte estas líneas en un dibujo. ¡Usa tu imaginación!}

\newcommand{\lblTituloPortada}{Aprendo a leer}
\newcommand{\lblSubtituloPortada}{Un capítulo nuevo cada día}
\newcommand{\lblPortadaNombre}{Nombre}
\newcommand{\lblPortadaCurso}{Curso}


% =========================================================================
% Paleta
% =========================================================================
% Cada actividad tiene un color saturado (borde + título) y uno claro
% (fondo). El color es un canal EXTRA, nunca el único: el título en
% negrita y el borde grueso ya distinguen la caja en blanco y negro -- lo
% que es lo que hace posible, sin cambiar una sola macro de las de abajo,
% una versión sin color de verdad (ver \bookcolor, fijado por main*.tex
% antes de % preamble.tex
% -----------------------------------------------------------------------
% "Aprendo a leer" -- cuaderno diario de lectura (español, A4).
% Todo el aparato tipográfico vive aquí: paquetes, colores, estilos de
% caja y las macros que usa cada página diaria. Requiere LuaLaTeX
% (fontspec) -- no pdflatex puro: la letra es Andika (SIL, OFL,
% fonts/andika/), pensada para primeros lectores -- "a"/"g" de un solo
% piso y "b"/"d"/"p"/"q" bien diferenciadas, a diferencia de Latin
% Modern Sans (ver notes/02-revision-y-plan.md, punto 5). fontspec ya
% cubre UTF-8 y la codificación de fuente por su cuenta -- no hacen
% falta fontenc ni inputenc, y de hecho estorban bajo LuaLaTeX.
% -----------------------------------------------------------------------

\usepackage{fontspec}
\setmainfont{Andika}[
  Path = fonts/andika/,
  Extension = .ttf,
  UprightFont = *-Regular,
  BoldFont = *-Bold,
  ItalicFont = *-Italic,
  BoldItalicFont = *-BoldItalic,
]
\setsansfont{Andika}[
  Path = fonts/andika/,
  Extension = .ttf,
  UprightFont = *-Regular,
  BoldFont = *-Bold,
  ItalicFont = *-Italic,
  BoldItalicFont = *-BoldItalic,
]

\usepackage[spanish]{babel}
% babel-spanish makes " an ACTIVE character (shorthand for <</>> and other
% combinations) -- a literal " typed in prose (dialogue, "quoted speech")
% collides with it and garbles silently, worst right before a capital
% letter ("E -> a stray accent, e.g. "Estoy" rendered as ".Ëtoy"). Plain
% \shorthandoff{"} here in the preamble does NOT work and fails silently
% -- babel re-activates its shorthands for the main language at
% \begin{document}, AFTER the preamble has already run, so it wins.
% \AtBeginDocument defers the command to run right after that, which is
% the only point it sticks -- verified against a minimal standalone file
% both ways before trusting either.
\AtBeginDocument{\shorthandoff{"}}

\usepackage[a4paper,top=18mm,bottom=16mm,left=22mm,right=18mm]{geometry}

\usepackage{xcolor}
\usepackage[most]{tcolorbox}
\usepackage{tabularx}
\usepackage{booktabs}
\usepackage{enumitem}
\usepackage{tikz}
\usepackage{multicol}

\pagestyle{empty}
\setlength{\parindent}{0pt}
\setlength{\parskip}{2pt}

% --- cadenas de texto (español; centralizadas, ver lang/es.tex) -------
% lang/es.tex
% -----------------------------------------------------------------------
% Todas las cadenas de texto que aparecen en el cuaderno, en un solo
% sitio. Cambiar una palabra aquí la cambia en las 260 páginas a la vez,
% sin tocar el contenido generado.
% -----------------------------------------------------------------------

\newcommand{\lblFecha}{Fecha}
\newcommand{\lblNotas}{Notas}
\newcommand{\lblDia}{Día}
\newcommand{\lblSemana}{Semana}
\newcommand{\lblTrimestre}{Trimestre}
\newcommand{\lblRespuesta}{Respuesta}

\newcommand{\lblInstruccionLectura}{Lee en voz alta, despacio, señalando cada palabra}

\newcommand{\lblDibuja}{Dibuja}
\newcommand{\lblCompleta}{Completa}
\newcommand{\lblCopia}{Copia}
\newcommand{\lblResponde}{Responde}
\newcommand{\lblRelaciona}{Relaciona}
\newcommand{\lblAdivina}{Adivina}
\newcommand{\lblCrea}{Crea}
\newcommand{\lblRepasa}{Repasa}

\newcommand{\lblInstruccionCopia}{Copia la frase de hoy, tres veces, una en cada línea.}
\newcommand{\lblInstruccionCompleta}{Convierte estas líneas en un dibujo. ¡Usa tu imaginación!}

\newcommand{\lblTituloPortada}{Aprendo a leer}
\newcommand{\lblSubtituloPortada}{Un capítulo nuevo cada día}
\newcommand{\lblPortadaNombre}{Nombre}
\newcommand{\lblPortadaCurso}{Curso}


% =========================================================================
% Paleta
% =========================================================================
% Cada actividad tiene un color saturado (borde + título) y uno claro
% (fondo). El color es un canal EXTRA, nunca el único: el título en
% negrita y el borde grueso ya distinguen la caja en blanco y negro -- lo
% que es lo que hace posible, sin cambiar una sola macro de las de abajo,
% una versión sin color de verdad (ver \bookcolor, fijado por main*.tex
% antes de % preamble.tex
% -----------------------------------------------------------------------
% "Aprendo a leer" -- cuaderno diario de lectura (español, A4).
% Todo el aparato tipográfico vive aquí: paquetes, colores, estilos de
% caja y las macros que usa cada página diaria. Requiere LuaLaTeX
% (fontspec) -- no pdflatex puro: la letra es Andika (SIL, OFL,
% fonts/andika/), pensada para primeros lectores -- "a"/"g" de un solo
% piso y "b"/"d"/"p"/"q" bien diferenciadas, a diferencia de Latin
% Modern Sans (ver notes/02-revision-y-plan.md, punto 5). fontspec ya
% cubre UTF-8 y la codificación de fuente por su cuenta -- no hacen
% falta fontenc ni inputenc, y de hecho estorban bajo LuaLaTeX.
% -----------------------------------------------------------------------

\usepackage{fontspec}
\setmainfont{Andika}[
  Path = fonts/andika/,
  Extension = .ttf,
  UprightFont = *-Regular,
  BoldFont = *-Bold,
  ItalicFont = *-Italic,
  BoldItalicFont = *-BoldItalic,
]
\setsansfont{Andika}[
  Path = fonts/andika/,
  Extension = .ttf,
  UprightFont = *-Regular,
  BoldFont = *-Bold,
  ItalicFont = *-Italic,
  BoldItalicFont = *-BoldItalic,
]

\usepackage[spanish]{babel}
% babel-spanish makes " an ACTIVE character (shorthand for <</>> and other
% combinations) -- a literal " typed in prose (dialogue, "quoted speech")
% collides with it and garbles silently, worst right before a capital
% letter ("E -> a stray accent, e.g. "Estoy" rendered as ".Ëtoy"). Plain
% \shorthandoff{"} here in the preamble does NOT work and fails silently
% -- babel re-activates its shorthands for the main language at
% \begin{document}, AFTER the preamble has already run, so it wins.
% \AtBeginDocument defers the command to run right after that, which is
% the only point it sticks -- verified against a minimal standalone file
% both ways before trusting either.
\AtBeginDocument{\shorthandoff{"}}

\usepackage[a4paper,top=18mm,bottom=16mm,left=22mm,right=18mm]{geometry}

\usepackage{xcolor}
\usepackage[most]{tcolorbox}
\usepackage{tabularx}
\usepackage{booktabs}
\usepackage{enumitem}
\usepackage{tikz}
\usepackage{multicol}

\pagestyle{empty}
\setlength{\parindent}{0pt}
\setlength{\parskip}{2pt}

% --- cadenas de texto (español; centralizadas, ver lang/es.tex) -------
% lang/es.tex
% -----------------------------------------------------------------------
% Todas las cadenas de texto que aparecen en el cuaderno, en un solo
% sitio. Cambiar una palabra aquí la cambia en las 260 páginas a la vez,
% sin tocar el contenido generado.
% -----------------------------------------------------------------------

\newcommand{\lblFecha}{Fecha}
\newcommand{\lblNotas}{Notas}
\newcommand{\lblDia}{Día}
\newcommand{\lblSemana}{Semana}
\newcommand{\lblTrimestre}{Trimestre}
\newcommand{\lblRespuesta}{Respuesta}

\newcommand{\lblInstruccionLectura}{Lee en voz alta, despacio, señalando cada palabra}

\newcommand{\lblDibuja}{Dibuja}
\newcommand{\lblCompleta}{Completa}
\newcommand{\lblCopia}{Copia}
\newcommand{\lblResponde}{Responde}
\newcommand{\lblRelaciona}{Relaciona}
\newcommand{\lblAdivina}{Adivina}
\newcommand{\lblCrea}{Crea}
\newcommand{\lblRepasa}{Repasa}

\newcommand{\lblInstruccionCopia}{Copia la frase de hoy, tres veces, una en cada línea.}
\newcommand{\lblInstruccionCompleta}{Convierte estas líneas en un dibujo. ¡Usa tu imaginación!}

\newcommand{\lblTituloPortada}{Aprendo a leer}
\newcommand{\lblSubtituloPortada}{Un capítulo nuevo cada día}
\newcommand{\lblPortadaNombre}{Nombre}
\newcommand{\lblPortadaCurso}{Curso}


% =========================================================================
% Paleta
% =========================================================================
% Cada actividad tiene un color saturado (borde + título) y uno claro
% (fondo). El color es un canal EXTRA, nunca el único: el título en
% negrita y el borde grueso ya distinguen la caja en blanco y negro -- lo
% que es lo que hace posible, sin cambiar una sola macro de las de abajo,
% una versión sin color de verdad (ver \bookcolor, fijado por main*.tex
% antes de % preamble.tex
% -----------------------------------------------------------------------
% "Aprendo a leer" -- cuaderno diario de lectura (español, A4).
% Todo el aparato tipográfico vive aquí: paquetes, colores, estilos de
% caja y las macros que usa cada página diaria. Requiere LuaLaTeX
% (fontspec) -- no pdflatex puro: la letra es Andika (SIL, OFL,
% fonts/andika/), pensada para primeros lectores -- "a"/"g" de un solo
% piso y "b"/"d"/"p"/"q" bien diferenciadas, a diferencia de Latin
% Modern Sans (ver notes/02-revision-y-plan.md, punto 5). fontspec ya
% cubre UTF-8 y la codificación de fuente por su cuenta -- no hacen
% falta fontenc ni inputenc, y de hecho estorban bajo LuaLaTeX.
% -----------------------------------------------------------------------

\usepackage{fontspec}
\setmainfont{Andika}[
  Path = fonts/andika/,
  Extension = .ttf,
  UprightFont = *-Regular,
  BoldFont = *-Bold,
  ItalicFont = *-Italic,
  BoldItalicFont = *-BoldItalic,
]
\setsansfont{Andika}[
  Path = fonts/andika/,
  Extension = .ttf,
  UprightFont = *-Regular,
  BoldFont = *-Bold,
  ItalicFont = *-Italic,
  BoldItalicFont = *-BoldItalic,
]

\usepackage[spanish]{babel}
% babel-spanish makes " an ACTIVE character (shorthand for <</>> and other
% combinations) -- a literal " typed in prose (dialogue, "quoted speech")
% collides with it and garbles silently, worst right before a capital
% letter ("E -> a stray accent, e.g. "Estoy" rendered as ".Ëtoy"). Plain
% \shorthandoff{"} here in the preamble does NOT work and fails silently
% -- babel re-activates its shorthands for the main language at
% \begin{document}, AFTER the preamble has already run, so it wins.
% \AtBeginDocument defers the command to run right after that, which is
% the only point it sticks -- verified against a minimal standalone file
% both ways before trusting either.
\AtBeginDocument{\shorthandoff{"}}

\usepackage[a4paper,top=18mm,bottom=16mm,left=22mm,right=18mm]{geometry}

\usepackage{xcolor}
\usepackage[most]{tcolorbox}
\usepackage{tabularx}
\usepackage{booktabs}
\usepackage{enumitem}
\usepackage{tikz}
\usepackage{multicol}

\pagestyle{empty}
\setlength{\parindent}{0pt}
\setlength{\parskip}{2pt}

% --- cadenas de texto (español; centralizadas, ver lang/es.tex) -------
\input{lang/es}

% =========================================================================
% Paleta
% =========================================================================
% Cada actividad tiene un color saturado (borde + título) y uno claro
% (fondo). El color es un canal EXTRA, nunca el único: el título en
% negrita y el borde grueso ya distinguen la caja en blanco y negro -- lo
% que es lo que hace posible, sin cambiar una sola macro de las de abajo,
% una versión sin color de verdad (ver \bookcolor, fijado por main*.tex
% antes de \input{preamble}): no se pierde ninguna información, solo tinta.
\makeatletter
\newif\ifbwbook
\def\aal@bookcolor@bw{bw}
\@ifundefined{bookcolor}{\bwbookfalse}{%
  \ifx\bookcolor\aal@bookcolor@bw \bwbooktrue \else \bwbookfalse \fi
}
\makeatother

\ifbwbook
  % Modo blanco y negro: mismos nombres de color, apuntan a negro/gris en
  % vez de a un tono por actividad -- ninguna caja de abajo sabe ni le
  % importa en qué modo está.
  \definecolor{colorLectura}{gray}{0}
  \definecolor{colorLecturaClaro}{gray}{0.95}
  \definecolor{colorDibuja}{gray}{0}
  \definecolor{colorDibujaClaro}{gray}{0.95}
  \definecolor{colorCompleta}{gray}{0}
  \definecolor{colorCompletaClaro}{gray}{0.95}
  \definecolor{colorResponde}{gray}{0}
  \definecolor{colorRespondeClaro}{gray}{0.95}
  \definecolor{colorRelaciona}{gray}{0}
  \definecolor{colorRelacionaClaro}{gray}{0.95}
  \definecolor{colorCrea}{gray}{0}
  \definecolor{colorCreaClaro}{gray}{0.95}
  \definecolor{colorGris}{gray}{0.38}
\else
  \definecolor{colorLectura}{HTML}{2E7D32}
  \definecolor{colorLecturaClaro}{HTML}{EAF6EC}
  \definecolor{colorDibuja}{HTML}{1565C0}
  \definecolor{colorDibujaClaro}{HTML}{E7F0FB}
  \definecolor{colorCompleta}{HTML}{C2185B}
  \definecolor{colorCompletaClaro}{HTML}{FBE9EF}
  \definecolor{colorResponde}{HTML}{6A1B9A}
  \definecolor{colorRespondeClaro}{HTML}{F2E9F7}
  \definecolor{colorRelaciona}{HTML}{E65100}
  \definecolor{colorRelacionaClaro}{HTML}{FDEEE2}
  \definecolor{colorCrea}{HTML}{AD1457}
  \definecolor{colorCreaClaro}{HTML}{FBE7EE}
  \definecolor{colorGris}{HTML}{616161}
\fi

% =========================================================================
% Cajas (tcolorbox)
% =========================================================================
\tcbset{
  cajabase/.style={
    enhanced,
    boxrule=1.1pt,
    arc=3mm,
    left=6mm, right=6mm, top=6mm, bottom=6mm,
    fonttitle=\bfseries\sffamily\large,
    coltitle=white,
  }
}

% El título es un argumento obligatorio (#1) en vez de un valor fijo:
% tools/gen_days.py elige qué instrucción pasar según el trimestre de
% cada día (ver PLANTILLA_DIA, \begin{cajaLectura}{...}). Probado a
% mano que \begin{cajaLectura}[title=...] (la sintaxis "de toda la
% vida" para sobreescribir una opción de tcolorbox) NO funciona aquí --
% \newtcolorbox sin un número de argumentos declarado no admite ese
% corchete, y el texto "[title=...]" acaba escrito literalmente dentro
% de la caja; declarar el título como argumento obligatorio es la forma
% que sí funciona.
% before upper bloquea la partición de palabras SOLO dentro de esta
% caja -- ver notes/02-revision-y-plan.md, punto 5: en cuanto el texto
% deja de estar centrado (trimestre 2 en adelante), LaTeX empieza a
% partir palabras a final de línea para justificar, y a los 6 años eso
% es ruido, no ayuda.
\newtcolorbox{cajaLectura}[1]{cajabase,
  colback=colorLecturaClaro, colframe=colorLectura, colbacktitle=colorLectura,
  title=#1,
  before upper={\hyphenpenalty=10000\exhyphenpenalty=10000}}

\newtcolorbox{cajaDibuja}{cajabase,
  colback=colorDibujaClaro, colframe=colorDibuja, colbacktitle=colorDibuja,
  title=\lblDibuja}

\newtcolorbox{cajaCompleta}{cajabase,
  colback=colorCompletaClaro, colframe=colorCompleta, colbacktitle=colorCompleta,
  title=\lblCompleta}

\newtcolorbox{cajaCopia}{cajabase,
  colback=colorRespondeClaro, colframe=colorResponde, colbacktitle=colorResponde,
  title=\lblCopia}

\newtcolorbox{cajaResponde}{cajabase,
  colback=colorRespondeClaro, colframe=colorResponde, colbacktitle=colorResponde,
  title=\lblResponde}

\newtcolorbox{cajaRelaciona}{cajabase,
  colback=colorRelacionaClaro, colframe=colorRelaciona, colbacktitle=colorRelaciona,
  title=\lblRelaciona}

\newtcolorbox{cajaAdivina}{cajabase,
  colback=colorRelacionaClaro, colframe=colorRelaciona, colbacktitle=colorRelaciona,
  title=\lblAdivina}

\newtcolorbox{cajaCrea}{cajabase,
  colback=colorCreaClaro, colframe=colorCrea, colbacktitle=colorCrea,
  title=\lblCrea}

\newtcolorbox{cajaRepasa}{cajabase,
  colback=colorCreaClaro, colframe=colorCrea, colbacktitle=colorCrea,
  title=\lblRepasa}

% Los cinco tipos nuevos del ciclo objetivo (ver
% notes/02-revision-y-plan.md, punto 3): reutilizan el color de la
% actividad más parecida en vez de añadir tonos nuevos a la paleta --
% rodea/verdadero_falso/ordena son variantes de "comprobar sin escribir
% mucho" (como Responde/Copia), busca es una tarea a nivel de palabra
% (como Completa), y relee es, literalmente, releer (como la propia
% caja de lectura).
\newtcolorbox{cajaRodea}{cajabase,
  colback=colorRespondeClaro, colframe=colorResponde, colbacktitle=colorResponde,
  title=\lblRodea}

\newtcolorbox{cajaVerdaderoFalso}{cajabase,
  colback=colorRespondeClaro, colframe=colorResponde, colbacktitle=colorResponde,
  title=\lblVerdaderoFalso}

\newtcolorbox{cajaBusca}{cajabase,
  colback=colorCompletaClaro, colframe=colorCompleta, colbacktitle=colorCompleta,
  title=\lblBusca}

\newtcolorbox{cajaOrdena}{cajabase,
  colback=colorRespondeClaro, colframe=colorResponde, colbacktitle=colorResponde,
  title=\lblOrdena}

\newtcolorbox{cajaRelee}{cajabase,
  colback=colorLecturaClaro, colframe=colorLectura, colbacktitle=colorLectura,
  title=\lblRelee}

% "Traza" reutiliza el color de "Copia"/"Responde" -- las dos son,
% igual que "traza", practicar a mano (letra o respuesta), no leer.
\newtcolorbox{cajaTraza}{cajabase,
  colback=colorRespondeClaro, colframe=colorResponde, colbacktitle=colorResponde,
  title=\lblTraza}

% Envoltorios de actividad: el color/título ya está fijado arriba; el
% contenido concreto (texto, tabla, dibujo, espacio en blanco) se escribe
% al llamar a la macro, así que cada día es libre de necesitar cosas
% distintas sin que la macro tenga que adivinarlas.
\newcommand{\actividadDibuja}[1]{\begin{cajaDibuja}#1\end{cajaDibuja}}
% A diferencia de las demás, la plantilla de "completa" (tools/gen_days.py)
% necesita meter una línea de instrucción a tamaño normal ANTES del dibujo
% escalado, así que aquí no se envuelve nada en \resizebox -- eso lo hace
% la propia plantilla, solo alrededor del \input{diagrams/...}.
\newcommand{\actividadCompleta}[1]{\begin{cajaCompleta}#1\end{cajaCompleta}}
\newcommand{\actividadCopia}[1]{\begin{cajaCopia}\large#1\end{cajaCopia}}
\newcommand{\actividadResponde}[1]{\begin{cajaResponde}\large#1\end{cajaResponde}}
\newcommand{\actividadRelaciona}[1]{\begin{cajaRelaciona}\large#1\end{cajaRelaciona}}
\newcommand{\actividadAdivina}[1]{\begin{cajaAdivina}#1\end{cajaAdivina}}
\newcommand{\actividadCrea}[1]{\begin{cajaCrea}#1\end{cajaCrea}}
\newcommand{\actividadRepasa}[1]{\begin{cajaRepasa}#1\end{cajaRepasa}}
\newcommand{\actividadRodea}[1]{\begin{cajaRodea}#1\end{cajaRodea}}
\newcommand{\actividadVerdaderoFalso}[1]{\begin{cajaVerdaderoFalso}#1\end{cajaVerdaderoFalso}}
\newcommand{\actividadBusca}[1]{\begin{cajaBusca}\large#1\end{cajaBusca}}
\newcommand{\actividadOrdena}[1]{\begin{cajaOrdena}#1\end{cajaOrdena}}
\newcommand{\actividadRelee}[1]{\begin{cajaRelee}\large#1\end{cajaRelee}}
\newcommand{\actividadTraza}[1]{\begin{cajaTraza}#1\end{cajaTraza}}

% Piezas reutilizables dentro de una actividad. El bloque de lectura de
% T1 deja mucho sitio en la página, así que el espacio en blanco por
% defecto es generoso -- que la caja de actividad tenga un peso visual
% parecido al de la caja de lectura, en vez de quedarse pequeña con
% media página vacía debajo.
\newcommand{\espacioDibujo}[1][11cm]{\vspace{#1}}
\newcommand{\lineaRespuesta}{\par\vspace{13mm}\noindent\rule{\linewidth}{0.4pt}}
\newcommand{\casilla}{\raisebox{-0.15ex}{\framebox[3.4mm][c]{\rule{0pt}{2.6mm}}}}
% Como \casilla, pero más grande -- para escribir un número (1, 2, 3) a
% mano dentro, en vez de solo marcar con una X. Usada por "ordena".
\newcommand{\casillaNumero}{\raisebox{-0.6ex}{\framebox[6mm][c]{\rule{0pt}{5mm}}}}
\newenvironment{listaRepaso}{%
  \begin{itemize}[label=\casilla, leftmargin=8mm, itemsep=2mm, topsep=2mm]
}{%
  \end{itemize}
}

% =========================================================================
% Clave de respuestas (backmatter/clave-respuestas.tex, generada por
% tools/gen_days.py en content/generated-clave.tex). Una entrada por
% cada día "adivina" o "relaciona" -- ver la revisión de un lector
% externo: sin esto, un adulto no puede comprobar una adivinanza o un
% relaciona sin resolverlos él mismo primero.
% =========================================================================
\newcommand{\claveEntrada}[3]{%
  \item \textbf{\lblDia\ #1} \textcolor{colorGris}{(#2)}: #3%
}

% =========================================================================
% Tamaño de la frase de lectura, según el trimestre.
% Trimestre 1: una frase, muy grande. Trimestre 4: cuatro frases, algo
% más pequeñas para que quepan en el mismo bloque de la página.
% =========================================================================
\newcommand{\diafuente}[1]{%
  \ifcase#1\relax
  \or \fontsize{28}{33}\selectfont
  \or \fontsize{22}{27}\selectfont
  \or \fontsize{18}{23}\selectfont
  \or \fontsize{15}{20}\selectfont
  \fi
}

% =========================================================================
% La página de un día.
%
%   \begin{diapagina}{numero}{semana}{trimestre}{tema}
%     \begin{cajaLectura}
%     \centering\diafuente{trimestre}
%     Texto de lectura (una o varias frases).
%     \end{cajaLectura}
%     \vspace{5mm}
%     \actividadX{...}
%   \end{diapagina}
%
% \begin{diapagina} pone la cabecera (Día/Semana/Trimestre, tema, Fecha,
% Leído, Notas) y \label{dia:N}; \end{diapagina} pone el pie y \newpage.
% Todo lo de en medio (la caja de lectura y la actividad) se escribe
% explícito, para que cada día pueda tener una actividad completamente
% distinta sin que la macro tenga que saber de antemano qué forma tiene.
%
% "Leído: [] sola [] con ayuda [] con dificultad" sustituye a un campo
% de Notas más grande -- ver notes/02-revision-y-plan.md, punto 5: tres
% casillas que un adulto marca en 2 segundos, cada día, terminan siendo
% un registro de progreso de todo el año sin ningún esfuerzo extra. El
% campo Notas se queda, más corto, para lo que no cabe en una casilla.
%
% La cabecera vive en su propia macro, \cabeceraDia, porque la comparten
% los dos cuadernos: este (diapagina) y el de primeras palabras
% (palabrapagina, ver preamble-palabras.tex) -- la misma cabecera en
% los dos, sin dos copias que puedan divergir.
% =========================================================================
\newcommand{\cabeceraDia}[4]{%
  \label{dia:#1}%
  \noindent{\LARGE\bfseries\lblDia~#1}\quad
  {\large(\lblSemana~#2 \textperiodcentered\ \lblTrimestre~#3)}\par
  {\normalsize\bfseries\color{colorGris}#4}\par
  \vspace{1mm}
  \textbf{\lblFecha:} \rule{22mm}{0.4pt}\hspace{4mm}%
  \textbf{\lblLeido:} \casilla\ \lblLeidoSola\quad
  \casilla\ \lblLeidoConAyuda\quad
  \casilla\ \lblLeidoConDificultad\hspace{4mm}%
  \textbf{\lblNotas:} \rule{20mm}{0.4pt}\par
  \vspace{3mm}
}
\newenvironment{diapagina}[4]{%
  \def\diapagina@actual{#1}%
  \cabeceraDia{#1}{#2}{#3}{#4}%
}{%
  \vfill
  \begin{center}\footnotesize\color{colorGris}\lblDia~\diapagina@actual~/ \totaldias\end{center}
  \newpage
}
): no se pierde ninguna información, solo tinta.
\makeatletter
\newif\ifbwbook
\def\aal@bookcolor@bw{bw}
\@ifundefined{bookcolor}{\bwbookfalse}{%
  \ifx\bookcolor\aal@bookcolor@bw \bwbooktrue \else \bwbookfalse \fi
}
\makeatother

\ifbwbook
  % Modo blanco y negro: mismos nombres de color, apuntan a negro/gris en
  % vez de a un tono por actividad -- ninguna caja de abajo sabe ni le
  % importa en qué modo está.
  \definecolor{colorLectura}{gray}{0}
  \definecolor{colorLecturaClaro}{gray}{0.95}
  \definecolor{colorDibuja}{gray}{0}
  \definecolor{colorDibujaClaro}{gray}{0.95}
  \definecolor{colorCompleta}{gray}{0}
  \definecolor{colorCompletaClaro}{gray}{0.95}
  \definecolor{colorResponde}{gray}{0}
  \definecolor{colorRespondeClaro}{gray}{0.95}
  \definecolor{colorRelaciona}{gray}{0}
  \definecolor{colorRelacionaClaro}{gray}{0.95}
  \definecolor{colorCrea}{gray}{0}
  \definecolor{colorCreaClaro}{gray}{0.95}
  \definecolor{colorGris}{gray}{0.38}
\else
  \definecolor{colorLectura}{HTML}{2E7D32}
  \definecolor{colorLecturaClaro}{HTML}{EAF6EC}
  \definecolor{colorDibuja}{HTML}{1565C0}
  \definecolor{colorDibujaClaro}{HTML}{E7F0FB}
  \definecolor{colorCompleta}{HTML}{C2185B}
  \definecolor{colorCompletaClaro}{HTML}{FBE9EF}
  \definecolor{colorResponde}{HTML}{6A1B9A}
  \definecolor{colorRespondeClaro}{HTML}{F2E9F7}
  \definecolor{colorRelaciona}{HTML}{E65100}
  \definecolor{colorRelacionaClaro}{HTML}{FDEEE2}
  \definecolor{colorCrea}{HTML}{AD1457}
  \definecolor{colorCreaClaro}{HTML}{FBE7EE}
  \definecolor{colorGris}{HTML}{616161}
\fi

% =========================================================================
% Cajas (tcolorbox)
% =========================================================================
\tcbset{
  cajabase/.style={
    enhanced,
    boxrule=1.1pt,
    arc=3mm,
    left=6mm, right=6mm, top=6mm, bottom=6mm,
    fonttitle=\bfseries\sffamily\large,
    coltitle=white,
  }
}

% El título es un argumento obligatorio (#1) en vez de un valor fijo:
% tools/gen_days.py elige qué instrucción pasar según el trimestre de
% cada día (ver PLANTILLA_DIA, \begin{cajaLectura}{...}). Probado a
% mano que \begin{cajaLectura}[title=...] (la sintaxis "de toda la
% vida" para sobreescribir una opción de tcolorbox) NO funciona aquí --
% \newtcolorbox sin un número de argumentos declarado no admite ese
% corchete, y el texto "[title=...]" acaba escrito literalmente dentro
% de la caja; declarar el título como argumento obligatorio es la forma
% que sí funciona.
% before upper bloquea la partición de palabras SOLO dentro de esta
% caja -- ver notes/02-revision-y-plan.md, punto 5: en cuanto el texto
% deja de estar centrado (trimestre 2 en adelante), LaTeX empieza a
% partir palabras a final de línea para justificar, y a los 6 años eso
% es ruido, no ayuda.
\newtcolorbox{cajaLectura}[1]{cajabase,
  colback=colorLecturaClaro, colframe=colorLectura, colbacktitle=colorLectura,
  title=#1,
  before upper={\hyphenpenalty=10000\exhyphenpenalty=10000}}

\newtcolorbox{cajaDibuja}{cajabase,
  colback=colorDibujaClaro, colframe=colorDibuja, colbacktitle=colorDibuja,
  title=\lblDibuja}

\newtcolorbox{cajaCompleta}{cajabase,
  colback=colorCompletaClaro, colframe=colorCompleta, colbacktitle=colorCompleta,
  title=\lblCompleta}

\newtcolorbox{cajaCopia}{cajabase,
  colback=colorRespondeClaro, colframe=colorResponde, colbacktitle=colorResponde,
  title=\lblCopia}

\newtcolorbox{cajaResponde}{cajabase,
  colback=colorRespondeClaro, colframe=colorResponde, colbacktitle=colorResponde,
  title=\lblResponde}

\newtcolorbox{cajaRelaciona}{cajabase,
  colback=colorRelacionaClaro, colframe=colorRelaciona, colbacktitle=colorRelaciona,
  title=\lblRelaciona}

\newtcolorbox{cajaAdivina}{cajabase,
  colback=colorRelacionaClaro, colframe=colorRelaciona, colbacktitle=colorRelaciona,
  title=\lblAdivina}

\newtcolorbox{cajaCrea}{cajabase,
  colback=colorCreaClaro, colframe=colorCrea, colbacktitle=colorCrea,
  title=\lblCrea}

\newtcolorbox{cajaRepasa}{cajabase,
  colback=colorCreaClaro, colframe=colorCrea, colbacktitle=colorCrea,
  title=\lblRepasa}

% Los cinco tipos nuevos del ciclo objetivo (ver
% notes/02-revision-y-plan.md, punto 3): reutilizan el color de la
% actividad más parecida en vez de añadir tonos nuevos a la paleta --
% rodea/verdadero_falso/ordena son variantes de "comprobar sin escribir
% mucho" (como Responde/Copia), busca es una tarea a nivel de palabra
% (como Completa), y relee es, literalmente, releer (como la propia
% caja de lectura).
\newtcolorbox{cajaRodea}{cajabase,
  colback=colorRespondeClaro, colframe=colorResponde, colbacktitle=colorResponde,
  title=\lblRodea}

\newtcolorbox{cajaVerdaderoFalso}{cajabase,
  colback=colorRespondeClaro, colframe=colorResponde, colbacktitle=colorResponde,
  title=\lblVerdaderoFalso}

\newtcolorbox{cajaBusca}{cajabase,
  colback=colorCompletaClaro, colframe=colorCompleta, colbacktitle=colorCompleta,
  title=\lblBusca}

\newtcolorbox{cajaOrdena}{cajabase,
  colback=colorRespondeClaro, colframe=colorResponde, colbacktitle=colorResponde,
  title=\lblOrdena}

\newtcolorbox{cajaRelee}{cajabase,
  colback=colorLecturaClaro, colframe=colorLectura, colbacktitle=colorLectura,
  title=\lblRelee}

% "Traza" reutiliza el color de "Copia"/"Responde" -- las dos son,
% igual que "traza", practicar a mano (letra o respuesta), no leer.
\newtcolorbox{cajaTraza}{cajabase,
  colback=colorRespondeClaro, colframe=colorResponde, colbacktitle=colorResponde,
  title=\lblTraza}

% Envoltorios de actividad: el color/título ya está fijado arriba; el
% contenido concreto (texto, tabla, dibujo, espacio en blanco) se escribe
% al llamar a la macro, así que cada día es libre de necesitar cosas
% distintas sin que la macro tenga que adivinarlas.
\newcommand{\actividadDibuja}[1]{\begin{cajaDibuja}#1\end{cajaDibuja}}
% A diferencia de las demás, la plantilla de "completa" (tools/gen_days.py)
% necesita meter una línea de instrucción a tamaño normal ANTES del dibujo
% escalado, así que aquí no se envuelve nada en \resizebox -- eso lo hace
% la propia plantilla, solo alrededor del \input{diagrams/...}.
\newcommand{\actividadCompleta}[1]{\begin{cajaCompleta}#1\end{cajaCompleta}}
\newcommand{\actividadCopia}[1]{\begin{cajaCopia}\large#1\end{cajaCopia}}
\newcommand{\actividadResponde}[1]{\begin{cajaResponde}\large#1\end{cajaResponde}}
\newcommand{\actividadRelaciona}[1]{\begin{cajaRelaciona}\large#1\end{cajaRelaciona}}
\newcommand{\actividadAdivina}[1]{\begin{cajaAdivina}#1\end{cajaAdivina}}
\newcommand{\actividadCrea}[1]{\begin{cajaCrea}#1\end{cajaCrea}}
\newcommand{\actividadRepasa}[1]{\begin{cajaRepasa}#1\end{cajaRepasa}}
\newcommand{\actividadRodea}[1]{\begin{cajaRodea}#1\end{cajaRodea}}
\newcommand{\actividadVerdaderoFalso}[1]{\begin{cajaVerdaderoFalso}#1\end{cajaVerdaderoFalso}}
\newcommand{\actividadBusca}[1]{\begin{cajaBusca}\large#1\end{cajaBusca}}
\newcommand{\actividadOrdena}[1]{\begin{cajaOrdena}#1\end{cajaOrdena}}
\newcommand{\actividadRelee}[1]{\begin{cajaRelee}\large#1\end{cajaRelee}}
\newcommand{\actividadTraza}[1]{\begin{cajaTraza}#1\end{cajaTraza}}

% Piezas reutilizables dentro de una actividad. El bloque de lectura de
% T1 deja mucho sitio en la página, así que el espacio en blanco por
% defecto es generoso -- que la caja de actividad tenga un peso visual
% parecido al de la caja de lectura, en vez de quedarse pequeña con
% media página vacía debajo.
\newcommand{\espacioDibujo}[1][11cm]{\vspace{#1}}
\newcommand{\lineaRespuesta}{\par\vspace{13mm}\noindent\rule{\linewidth}{0.4pt}}
\newcommand{\casilla}{\raisebox{-0.15ex}{\framebox[3.4mm][c]{\rule{0pt}{2.6mm}}}}
% Como \casilla, pero más grande -- para escribir un número (1, 2, 3) a
% mano dentro, en vez de solo marcar con una X. Usada por "ordena".
\newcommand{\casillaNumero}{\raisebox{-0.6ex}{\framebox[6mm][c]{\rule{0pt}{5mm}}}}
\newenvironment{listaRepaso}{%
  \begin{itemize}[label=\casilla, leftmargin=8mm, itemsep=2mm, topsep=2mm]
}{%
  \end{itemize}
}

% =========================================================================
% Clave de respuestas (backmatter/clave-respuestas.tex, generada por
% tools/gen_days.py en content/generated-clave.tex). Una entrada por
% cada día "adivina" o "relaciona" -- ver la revisión de un lector
% externo: sin esto, un adulto no puede comprobar una adivinanza o un
% relaciona sin resolverlos él mismo primero.
% =========================================================================
\newcommand{\claveEntrada}[3]{%
  \item \textbf{\lblDia\ #1} \textcolor{colorGris}{(#2)}: #3%
}

% =========================================================================
% Tamaño de la frase de lectura, según el trimestre.
% Trimestre 1: una frase, muy grande. Trimestre 4: cuatro frases, algo
% más pequeñas para que quepan en el mismo bloque de la página.
% =========================================================================
\newcommand{\diafuente}[1]{%
  \ifcase#1\relax
  \or \fontsize{28}{33}\selectfont
  \or \fontsize{22}{27}\selectfont
  \or \fontsize{18}{23}\selectfont
  \or \fontsize{15}{20}\selectfont
  \fi
}

% =========================================================================
% La página de un día.
%
%   \begin{diapagina}{numero}{semana}{trimestre}{tema}
%     \begin{cajaLectura}
%     \centering\diafuente{trimestre}
%     Texto de lectura (una o varias frases).
%     \end{cajaLectura}
%     \vspace{5mm}
%     \actividadX{...}
%   \end{diapagina}
%
% \begin{diapagina} pone la cabecera (Día/Semana/Trimestre, tema, Fecha,
% Leído, Notas) y \label{dia:N}; \end{diapagina} pone el pie y \newpage.
% Todo lo de en medio (la caja de lectura y la actividad) se escribe
% explícito, para que cada día pueda tener una actividad completamente
% distinta sin que la macro tenga que saber de antemano qué forma tiene.
%
% "Leído: [] sola [] con ayuda [] con dificultad" sustituye a un campo
% de Notas más grande -- ver notes/02-revision-y-plan.md, punto 5: tres
% casillas que un adulto marca en 2 segundos, cada día, terminan siendo
% un registro de progreso de todo el año sin ningún esfuerzo extra. El
% campo Notas se queda, más corto, para lo que no cabe en una casilla.
%
% La cabecera vive en su propia macro, \cabeceraDia, porque la comparten
% los dos cuadernos: este (diapagina) y el de primeras palabras
% (palabrapagina, ver preamble-palabras.tex) -- la misma cabecera en
% los dos, sin dos copias que puedan divergir.
% =========================================================================
\newcommand{\cabeceraDia}[4]{%
  \label{dia:#1}%
  \noindent{\LARGE\bfseries\lblDia~#1}\quad
  {\large(\lblSemana~#2 \textperiodcentered\ \lblTrimestre~#3)}\par
  {\normalsize\bfseries\color{colorGris}#4}\par
  \vspace{1mm}
  \textbf{\lblFecha:} \rule{22mm}{0.4pt}\hspace{4mm}%
  \textbf{\lblLeido:} \casilla\ \lblLeidoSola\quad
  \casilla\ \lblLeidoConAyuda\quad
  \casilla\ \lblLeidoConDificultad\hspace{4mm}%
  \textbf{\lblNotas:} \rule{20mm}{0.4pt}\par
  \vspace{3mm}
}
\newenvironment{diapagina}[4]{%
  \def\diapagina@actual{#1}%
  \cabeceraDia{#1}{#2}{#3}{#4}%
}{%
  \vfill
  \begin{center}\footnotesize\color{colorGris}\lblDia~\diapagina@actual~/ \totaldias\end{center}
  \newpage
}
): no se pierde ninguna información, solo tinta.
\makeatletter
\newif\ifbwbook
\def\aal@bookcolor@bw{bw}
\@ifundefined{bookcolor}{\bwbookfalse}{%
  \ifx\bookcolor\aal@bookcolor@bw \bwbooktrue \else \bwbookfalse \fi
}
\makeatother

\ifbwbook
  % Modo blanco y negro: mismos nombres de color, apuntan a negro/gris en
  % vez de a un tono por actividad -- ninguna caja de abajo sabe ni le
  % importa en qué modo está.
  \definecolor{colorLectura}{gray}{0}
  \definecolor{colorLecturaClaro}{gray}{0.95}
  \definecolor{colorDibuja}{gray}{0}
  \definecolor{colorDibujaClaro}{gray}{0.95}
  \definecolor{colorCompleta}{gray}{0}
  \definecolor{colorCompletaClaro}{gray}{0.95}
  \definecolor{colorResponde}{gray}{0}
  \definecolor{colorRespondeClaro}{gray}{0.95}
  \definecolor{colorRelaciona}{gray}{0}
  \definecolor{colorRelacionaClaro}{gray}{0.95}
  \definecolor{colorCrea}{gray}{0}
  \definecolor{colorCreaClaro}{gray}{0.95}
  \definecolor{colorGris}{gray}{0.38}
\else
  \definecolor{colorLectura}{HTML}{2E7D32}
  \definecolor{colorLecturaClaro}{HTML}{EAF6EC}
  \definecolor{colorDibuja}{HTML}{1565C0}
  \definecolor{colorDibujaClaro}{HTML}{E7F0FB}
  \definecolor{colorCompleta}{HTML}{C2185B}
  \definecolor{colorCompletaClaro}{HTML}{FBE9EF}
  \definecolor{colorResponde}{HTML}{6A1B9A}
  \definecolor{colorRespondeClaro}{HTML}{F2E9F7}
  \definecolor{colorRelaciona}{HTML}{E65100}
  \definecolor{colorRelacionaClaro}{HTML}{FDEEE2}
  \definecolor{colorCrea}{HTML}{AD1457}
  \definecolor{colorCreaClaro}{HTML}{FBE7EE}
  \definecolor{colorGris}{HTML}{616161}
\fi

% =========================================================================
% Cajas (tcolorbox)
% =========================================================================
\tcbset{
  cajabase/.style={
    enhanced,
    boxrule=1.1pt,
    arc=3mm,
    left=6mm, right=6mm, top=6mm, bottom=6mm,
    fonttitle=\bfseries\sffamily\large,
    coltitle=white,
  }
}

% El título es un argumento obligatorio (#1) en vez de un valor fijo:
% tools/gen_days.py elige qué instrucción pasar según el trimestre de
% cada día (ver PLANTILLA_DIA, \begin{cajaLectura}{...}). Probado a
% mano que \begin{cajaLectura}[title=...] (la sintaxis "de toda la
% vida" para sobreescribir una opción de tcolorbox) NO funciona aquí --
% \newtcolorbox sin un número de argumentos declarado no admite ese
% corchete, y el texto "[title=...]" acaba escrito literalmente dentro
% de la caja; declarar el título como argumento obligatorio es la forma
% que sí funciona.
% before upper bloquea la partición de palabras SOLO dentro de esta
% caja -- ver notes/02-revision-y-plan.md, punto 5: en cuanto el texto
% deja de estar centrado (trimestre 2 en adelante), LaTeX empieza a
% partir palabras a final de línea para justificar, y a los 6 años eso
% es ruido, no ayuda.
\newtcolorbox{cajaLectura}[1]{cajabase,
  colback=colorLecturaClaro, colframe=colorLectura, colbacktitle=colorLectura,
  title=#1,
  before upper={\hyphenpenalty=10000\exhyphenpenalty=10000}}

\newtcolorbox{cajaDibuja}{cajabase,
  colback=colorDibujaClaro, colframe=colorDibuja, colbacktitle=colorDibuja,
  title=\lblDibuja}

\newtcolorbox{cajaCompleta}{cajabase,
  colback=colorCompletaClaro, colframe=colorCompleta, colbacktitle=colorCompleta,
  title=\lblCompleta}

\newtcolorbox{cajaCopia}{cajabase,
  colback=colorRespondeClaro, colframe=colorResponde, colbacktitle=colorResponde,
  title=\lblCopia}

\newtcolorbox{cajaResponde}{cajabase,
  colback=colorRespondeClaro, colframe=colorResponde, colbacktitle=colorResponde,
  title=\lblResponde}

\newtcolorbox{cajaRelaciona}{cajabase,
  colback=colorRelacionaClaro, colframe=colorRelaciona, colbacktitle=colorRelaciona,
  title=\lblRelaciona}

\newtcolorbox{cajaAdivina}{cajabase,
  colback=colorRelacionaClaro, colframe=colorRelaciona, colbacktitle=colorRelaciona,
  title=\lblAdivina}

\newtcolorbox{cajaCrea}{cajabase,
  colback=colorCreaClaro, colframe=colorCrea, colbacktitle=colorCrea,
  title=\lblCrea}

\newtcolorbox{cajaRepasa}{cajabase,
  colback=colorCreaClaro, colframe=colorCrea, colbacktitle=colorCrea,
  title=\lblRepasa}

% Los cinco tipos nuevos del ciclo objetivo (ver
% notes/02-revision-y-plan.md, punto 3): reutilizan el color de la
% actividad más parecida en vez de añadir tonos nuevos a la paleta --
% rodea/verdadero_falso/ordena son variantes de "comprobar sin escribir
% mucho" (como Responde/Copia), busca es una tarea a nivel de palabra
% (como Completa), y relee es, literalmente, releer (como la propia
% caja de lectura).
\newtcolorbox{cajaRodea}{cajabase,
  colback=colorRespondeClaro, colframe=colorResponde, colbacktitle=colorResponde,
  title=\lblRodea}

\newtcolorbox{cajaVerdaderoFalso}{cajabase,
  colback=colorRespondeClaro, colframe=colorResponde, colbacktitle=colorResponde,
  title=\lblVerdaderoFalso}

\newtcolorbox{cajaBusca}{cajabase,
  colback=colorCompletaClaro, colframe=colorCompleta, colbacktitle=colorCompleta,
  title=\lblBusca}

\newtcolorbox{cajaOrdena}{cajabase,
  colback=colorRespondeClaro, colframe=colorResponde, colbacktitle=colorResponde,
  title=\lblOrdena}

\newtcolorbox{cajaRelee}{cajabase,
  colback=colorLecturaClaro, colframe=colorLectura, colbacktitle=colorLectura,
  title=\lblRelee}

% "Traza" reutiliza el color de "Copia"/"Responde" -- las dos son,
% igual que "traza", practicar a mano (letra o respuesta), no leer.
\newtcolorbox{cajaTraza}{cajabase,
  colback=colorRespondeClaro, colframe=colorResponde, colbacktitle=colorResponde,
  title=\lblTraza}

% Envoltorios de actividad: el color/título ya está fijado arriba; el
% contenido concreto (texto, tabla, dibujo, espacio en blanco) se escribe
% al llamar a la macro, así que cada día es libre de necesitar cosas
% distintas sin que la macro tenga que adivinarlas.
\newcommand{\actividadDibuja}[1]{\begin{cajaDibuja}#1\end{cajaDibuja}}
% A diferencia de las demás, la plantilla de "completa" (tools/gen_days.py)
% necesita meter una línea de instrucción a tamaño normal ANTES del dibujo
% escalado, así que aquí no se envuelve nada en \resizebox -- eso lo hace
% la propia plantilla, solo alrededor del \input{diagrams/...}.
\newcommand{\actividadCompleta}[1]{\begin{cajaCompleta}#1\end{cajaCompleta}}
\newcommand{\actividadCopia}[1]{\begin{cajaCopia}\large#1\end{cajaCopia}}
\newcommand{\actividadResponde}[1]{\begin{cajaResponde}\large#1\end{cajaResponde}}
\newcommand{\actividadRelaciona}[1]{\begin{cajaRelaciona}\large#1\end{cajaRelaciona}}
\newcommand{\actividadAdivina}[1]{\begin{cajaAdivina}#1\end{cajaAdivina}}
\newcommand{\actividadCrea}[1]{\begin{cajaCrea}#1\end{cajaCrea}}
\newcommand{\actividadRepasa}[1]{\begin{cajaRepasa}#1\end{cajaRepasa}}
\newcommand{\actividadRodea}[1]{\begin{cajaRodea}#1\end{cajaRodea}}
\newcommand{\actividadVerdaderoFalso}[1]{\begin{cajaVerdaderoFalso}#1\end{cajaVerdaderoFalso}}
\newcommand{\actividadBusca}[1]{\begin{cajaBusca}\large#1\end{cajaBusca}}
\newcommand{\actividadOrdena}[1]{\begin{cajaOrdena}#1\end{cajaOrdena}}
\newcommand{\actividadRelee}[1]{\begin{cajaRelee}\large#1\end{cajaRelee}}
\newcommand{\actividadTraza}[1]{\begin{cajaTraza}#1\end{cajaTraza}}

% Piezas reutilizables dentro de una actividad. El bloque de lectura de
% T1 deja mucho sitio en la página, así que el espacio en blanco por
% defecto es generoso -- que la caja de actividad tenga un peso visual
% parecido al de la caja de lectura, en vez de quedarse pequeña con
% media página vacía debajo.
\newcommand{\espacioDibujo}[1][11cm]{\vspace{#1}}
\newcommand{\lineaRespuesta}{\par\vspace{13mm}\noindent\rule{\linewidth}{0.4pt}}
\newcommand{\casilla}{\raisebox{-0.15ex}{\framebox[3.4mm][c]{\rule{0pt}{2.6mm}}}}
% Como \casilla, pero más grande -- para escribir un número (1, 2, 3) a
% mano dentro, en vez de solo marcar con una X. Usada por "ordena".
\newcommand{\casillaNumero}{\raisebox{-0.6ex}{\framebox[6mm][c]{\rule{0pt}{5mm}}}}
\newenvironment{listaRepaso}{%
  \begin{itemize}[label=\casilla, leftmargin=8mm, itemsep=2mm, topsep=2mm]
}{%
  \end{itemize}
}

% =========================================================================
% Clave de respuestas (backmatter/clave-respuestas.tex, generada por
% tools/gen_days.py en content/generated-clave.tex). Una entrada por
% cada día "adivina" o "relaciona" -- ver la revisión de un lector
% externo: sin esto, un adulto no puede comprobar una adivinanza o un
% relaciona sin resolverlos él mismo primero.
% =========================================================================
\newcommand{\claveEntrada}[3]{%
  \item \textbf{\lblDia\ #1} \textcolor{colorGris}{(#2)}: #3%
}

% =========================================================================
% Tamaño de la frase de lectura, según el trimestre.
% Trimestre 1: una frase, muy grande. Trimestre 4: cuatro frases, algo
% más pequeñas para que quepan en el mismo bloque de la página.
% =========================================================================
\newcommand{\diafuente}[1]{%
  \ifcase#1\relax
  \or \fontsize{28}{33}\selectfont
  \or \fontsize{22}{27}\selectfont
  \or \fontsize{18}{23}\selectfont
  \or \fontsize{15}{20}\selectfont
  \fi
}

% =========================================================================
% La página de un día.
%
%   \begin{diapagina}{numero}{semana}{trimestre}{tema}
%     \begin{cajaLectura}
%     \centering\diafuente{trimestre}
%     Texto de lectura (una o varias frases).
%     \end{cajaLectura}
%     \vspace{5mm}
%     \actividadX{...}
%   \end{diapagina}
%
% \begin{diapagina} pone la cabecera (Día/Semana/Trimestre, tema, Fecha,
% Leído, Notas) y \label{dia:N}; \end{diapagina} pone el pie y \newpage.
% Todo lo de en medio (la caja de lectura y la actividad) se escribe
% explícito, para que cada día pueda tener una actividad completamente
% distinta sin que la macro tenga que saber de antemano qué forma tiene.
%
% "Leído: [] sola [] con ayuda [] con dificultad" sustituye a un campo
% de Notas más grande -- ver notes/02-revision-y-plan.md, punto 5: tres
% casillas que un adulto marca en 2 segundos, cada día, terminan siendo
% un registro de progreso de todo el año sin ningún esfuerzo extra. El
% campo Notas se queda, más corto, para lo que no cabe en una casilla.
%
% La cabecera vive en su propia macro, \cabeceraDia, porque la comparten
% los dos cuadernos: este (diapagina) y el de primeras palabras
% (palabrapagina, ver preamble-palabras.tex) -- la misma cabecera en
% los dos, sin dos copias que puedan divergir.
% =========================================================================
\newcommand{\cabeceraDia}[4]{%
  \label{dia:#1}%
  \noindent{\LARGE\bfseries\lblDia~#1}\quad
  {\large(\lblSemana~#2 \textperiodcentered\ \lblTrimestre~#3)}\par
  {\normalsize\bfseries\color{colorGris}#4}\par
  \vspace{1mm}
  \textbf{\lblFecha:} \rule{22mm}{0.4pt}\hspace{4mm}%
  \textbf{\lblLeido:} \casilla\ \lblLeidoSola\quad
  \casilla\ \lblLeidoConAyuda\quad
  \casilla\ \lblLeidoConDificultad\hspace{4mm}%
  \textbf{\lblNotas:} \rule{20mm}{0.4pt}\par
  \vspace{3mm}
}
\newenvironment{diapagina}[4]{%
  \def\diapagina@actual{#1}%
  \cabeceraDia{#1}{#2}{#3}{#4}%
}{%
  \vfill
  \begin{center}\footnotesize\color{colorGris}\lblDia~\diapagina@actual~/ \totaldias\end{center}
  \newpage
}
): no se pierde ninguna información, solo tinta.
\makeatletter
\newif\ifbwbook
\def\aal@bookcolor@bw{bw}
\@ifundefined{bookcolor}{\bwbookfalse}{%
  \ifx\bookcolor\aal@bookcolor@bw \bwbooktrue \else \bwbookfalse \fi
}
\makeatother

\ifbwbook
  % Modo blanco y negro: mismos nombres de color, apuntan a negro/gris en
  % vez de a un tono por actividad -- ninguna caja de abajo sabe ni le
  % importa en qué modo está.
  \definecolor{colorLectura}{gray}{0}
  \definecolor{colorLecturaClaro}{gray}{0.95}
  \definecolor{colorDibuja}{gray}{0}
  \definecolor{colorDibujaClaro}{gray}{0.95}
  \definecolor{colorCompleta}{gray}{0}
  \definecolor{colorCompletaClaro}{gray}{0.95}
  \definecolor{colorResponde}{gray}{0}
  \definecolor{colorRespondeClaro}{gray}{0.95}
  \definecolor{colorRelaciona}{gray}{0}
  \definecolor{colorRelacionaClaro}{gray}{0.95}
  \definecolor{colorCrea}{gray}{0}
  \definecolor{colorCreaClaro}{gray}{0.95}
  \definecolor{colorGris}{gray}{0.38}
\else
  \definecolor{colorLectura}{HTML}{2E7D32}
  \definecolor{colorLecturaClaro}{HTML}{EAF6EC}
  \definecolor{colorDibuja}{HTML}{1565C0}
  \definecolor{colorDibujaClaro}{HTML}{E7F0FB}
  \definecolor{colorCompleta}{HTML}{C2185B}
  \definecolor{colorCompletaClaro}{HTML}{FBE9EF}
  \definecolor{colorResponde}{HTML}{6A1B9A}
  \definecolor{colorRespondeClaro}{HTML}{F2E9F7}
  \definecolor{colorRelaciona}{HTML}{E65100}
  \definecolor{colorRelacionaClaro}{HTML}{FDEEE2}
  \definecolor{colorCrea}{HTML}{AD1457}
  \definecolor{colorCreaClaro}{HTML}{FBE7EE}
  \definecolor{colorGris}{HTML}{616161}
\fi

% =========================================================================
% Cajas (tcolorbox)
% =========================================================================
\tcbset{
  cajabase/.style={
    enhanced,
    boxrule=1.1pt,
    arc=3mm,
    left=6mm, right=6mm, top=6mm, bottom=6mm,
    fonttitle=\bfseries\sffamily\large,
    coltitle=white,
  }
}

% El título es un argumento obligatorio (#1) en vez de un valor fijo:
% tools/gen_days.py elige qué instrucción pasar según el trimestre de
% cada día (ver PLANTILLA_DIA, \begin{cajaLectura}{...}). Probado a
% mano que \begin{cajaLectura}[title=...] (la sintaxis "de toda la
% vida" para sobreescribir una opción de tcolorbox) NO funciona aquí --
% \newtcolorbox sin un número de argumentos declarado no admite ese
% corchete, y el texto "[title=...]" acaba escrito literalmente dentro
% de la caja; declarar el título como argumento obligatorio es la forma
% que sí funciona.
% before upper bloquea la partición de palabras SOLO dentro de esta
% caja -- ver notes/02-revision-y-plan.md, punto 5: en cuanto el texto
% deja de estar centrado (trimestre 2 en adelante), LaTeX empieza a
% partir palabras a final de línea para justificar, y a los 6 años eso
% es ruido, no ayuda.
\newtcolorbox{cajaLectura}[1]{cajabase,
  colback=colorLecturaClaro, colframe=colorLectura, colbacktitle=colorLectura,
  title=#1,
  before upper={\hyphenpenalty=10000\exhyphenpenalty=10000}}

\newtcolorbox{cajaDibuja}{cajabase,
  colback=colorDibujaClaro, colframe=colorDibuja, colbacktitle=colorDibuja,
  title=\lblDibuja}

\newtcolorbox{cajaCompleta}{cajabase,
  colback=colorCompletaClaro, colframe=colorCompleta, colbacktitle=colorCompleta,
  title=\lblCompleta}

\newtcolorbox{cajaCopia}{cajabase,
  colback=colorRespondeClaro, colframe=colorResponde, colbacktitle=colorResponde,
  title=\lblCopia}

\newtcolorbox{cajaResponde}{cajabase,
  colback=colorRespondeClaro, colframe=colorResponde, colbacktitle=colorResponde,
  title=\lblResponde}

\newtcolorbox{cajaRelaciona}{cajabase,
  colback=colorRelacionaClaro, colframe=colorRelaciona, colbacktitle=colorRelaciona,
  title=\lblRelaciona}

\newtcolorbox{cajaAdivina}{cajabase,
  colback=colorRelacionaClaro, colframe=colorRelaciona, colbacktitle=colorRelaciona,
  title=\lblAdivina}

\newtcolorbox{cajaCrea}{cajabase,
  colback=colorCreaClaro, colframe=colorCrea, colbacktitle=colorCrea,
  title=\lblCrea}

\newtcolorbox{cajaRepasa}{cajabase,
  colback=colorCreaClaro, colframe=colorCrea, colbacktitle=colorCrea,
  title=\lblRepasa}

% Los cinco tipos nuevos del ciclo objetivo (ver
% notes/02-revision-y-plan.md, punto 3): reutilizan el color de la
% actividad más parecida en vez de añadir tonos nuevos a la paleta --
% rodea/verdadero_falso/ordena son variantes de "comprobar sin escribir
% mucho" (como Responde/Copia), busca es una tarea a nivel de palabra
% (como Completa), y relee es, literalmente, releer (como la propia
% caja de lectura).
\newtcolorbox{cajaRodea}{cajabase,
  colback=colorRespondeClaro, colframe=colorResponde, colbacktitle=colorResponde,
  title=\lblRodea}

\newtcolorbox{cajaVerdaderoFalso}{cajabase,
  colback=colorRespondeClaro, colframe=colorResponde, colbacktitle=colorResponde,
  title=\lblVerdaderoFalso}

\newtcolorbox{cajaBusca}{cajabase,
  colback=colorCompletaClaro, colframe=colorCompleta, colbacktitle=colorCompleta,
  title=\lblBusca}

\newtcolorbox{cajaOrdena}{cajabase,
  colback=colorRespondeClaro, colframe=colorResponde, colbacktitle=colorResponde,
  title=\lblOrdena}

\newtcolorbox{cajaRelee}{cajabase,
  colback=colorLecturaClaro, colframe=colorLectura, colbacktitle=colorLectura,
  title=\lblRelee}

% "Traza" reutiliza el color de "Copia"/"Responde" -- las dos son,
% igual que "traza", practicar a mano (letra o respuesta), no leer.
\newtcolorbox{cajaTraza}{cajabase,
  colback=colorRespondeClaro, colframe=colorResponde, colbacktitle=colorResponde,
  title=\lblTraza}

% Envoltorios de actividad: el color/título ya está fijado arriba; el
% contenido concreto (texto, tabla, dibujo, espacio en blanco) se escribe
% al llamar a la macro, así que cada día es libre de necesitar cosas
% distintas sin que la macro tenga que adivinarlas.
\newcommand{\actividadDibuja}[1]{\begin{cajaDibuja}#1\end{cajaDibuja}}
% A diferencia de las demás, la plantilla de "completa" (tools/gen_days.py)
% necesita meter una línea de instrucción a tamaño normal ANTES del dibujo
% escalado, así que aquí no se envuelve nada en \resizebox -- eso lo hace
% la propia plantilla, solo alrededor del \input{diagrams/...}.
\newcommand{\actividadCompleta}[1]{\begin{cajaCompleta}#1\end{cajaCompleta}}
\newcommand{\actividadCopia}[1]{\begin{cajaCopia}\large#1\end{cajaCopia}}
\newcommand{\actividadResponde}[1]{\begin{cajaResponde}\large#1\end{cajaResponde}}
\newcommand{\actividadRelaciona}[1]{\begin{cajaRelaciona}\large#1\end{cajaRelaciona}}
\newcommand{\actividadAdivina}[1]{\begin{cajaAdivina}#1\end{cajaAdivina}}
\newcommand{\actividadCrea}[1]{\begin{cajaCrea}#1\end{cajaCrea}}
\newcommand{\actividadRepasa}[1]{\begin{cajaRepasa}#1\end{cajaRepasa}}
\newcommand{\actividadRodea}[1]{\begin{cajaRodea}#1\end{cajaRodea}}
\newcommand{\actividadVerdaderoFalso}[1]{\begin{cajaVerdaderoFalso}#1\end{cajaVerdaderoFalso}}
\newcommand{\actividadBusca}[1]{\begin{cajaBusca}\large#1\end{cajaBusca}}
\newcommand{\actividadOrdena}[1]{\begin{cajaOrdena}#1\end{cajaOrdena}}
\newcommand{\actividadRelee}[1]{\begin{cajaRelee}\large#1\end{cajaRelee}}
\newcommand{\actividadTraza}[1]{\begin{cajaTraza}#1\end{cajaTraza}}

% Piezas reutilizables dentro de una actividad. El bloque de lectura de
% T1 deja mucho sitio en la página, así que el espacio en blanco por
% defecto es generoso -- que la caja de actividad tenga un peso visual
% parecido al de la caja de lectura, en vez de quedarse pequeña con
% media página vacía debajo.
\newcommand{\espacioDibujo}[1][11cm]{\vspace{#1}}
\newcommand{\lineaRespuesta}{\par\vspace{13mm}\noindent\rule{\linewidth}{0.4pt}}
\newcommand{\casilla}{\raisebox{-0.15ex}{\framebox[3.4mm][c]{\rule{0pt}{2.6mm}}}}
% Como \casilla, pero más grande -- para escribir un número (1, 2, 3) a
% mano dentro, en vez de solo marcar con una X. Usada por "ordena".
\newcommand{\casillaNumero}{\raisebox{-0.6ex}{\framebox[6mm][c]{\rule{0pt}{5mm}}}}
\newenvironment{listaRepaso}{%
  \begin{itemize}[label=\casilla, leftmargin=8mm, itemsep=2mm, topsep=2mm]
}{%
  \end{itemize}
}

% =========================================================================
% Clave de respuestas (backmatter/clave-respuestas.tex, generada por
% tools/gen_days.py en content/generated-clave.tex). Una entrada por
% cada día "adivina" o "relaciona" -- ver la revisión de un lector
% externo: sin esto, un adulto no puede comprobar una adivinanza o un
% relaciona sin resolverlos él mismo primero.
% =========================================================================
\newcommand{\claveEntrada}[3]{%
  \item \textbf{\lblDia\ #1} \textcolor{colorGris}{(#2)}: #3%
}

% =========================================================================
% Tamaño de la frase de lectura, según el trimestre.
% Trimestre 1: una frase, muy grande. Trimestre 4: cuatro frases, algo
% más pequeñas para que quepan en el mismo bloque de la página.
% =========================================================================
\newcommand{\diafuente}[1]{%
  \ifcase#1\relax
  \or \fontsize{28}{33}\selectfont
  \or \fontsize{22}{27}\selectfont
  \or \fontsize{18}{23}\selectfont
  \or \fontsize{15}{20}\selectfont
  \fi
}

% =========================================================================
% La página de un día.
%
%   \begin{diapagina}{numero}{semana}{trimestre}{tema}
%     \begin{cajaLectura}
%     \centering\diafuente{trimestre}
%     Texto de lectura (una o varias frases).
%     \end{cajaLectura}
%     \vspace{5mm}
%     \actividadX{...}
%   \end{diapagina}
%
% \begin{diapagina} pone la cabecera (Día/Semana/Trimestre, tema, Fecha,
% Leído, Notas) y \label{dia:N}; \end{diapagina} pone el pie y \newpage.
% Todo lo de en medio (la caja de lectura y la actividad) se escribe
% explícito, para que cada día pueda tener una actividad completamente
% distinta sin que la macro tenga que saber de antemano qué forma tiene.
%
% "Leído: [] sola [] con ayuda [] con dificultad" sustituye a un campo
% de Notas más grande -- ver notes/02-revision-y-plan.md, punto 5: tres
% casillas que un adulto marca en 2 segundos, cada día, terminan siendo
% un registro de progreso de todo el año sin ningún esfuerzo extra. El
% campo Notas se queda, más corto, para lo que no cabe en una casilla.
%
% La cabecera vive en su propia macro, \cabeceraDia, porque la comparten
% los dos cuadernos: este (diapagina) y el de primeras palabras
% (palabrapagina, ver preamble-palabras.tex) -- la misma cabecera en
% los dos, sin dos copias que puedan divergir.
% =========================================================================
\newcommand{\cabeceraDia}[4]{%
  \label{dia:#1}%
  \noindent{\LARGE\bfseries\lblDia~#1}\quad
  {\large(\lblSemana~#2 \textperiodcentered\ \lblTrimestre~#3)}\par
  {\normalsize\bfseries\color{colorGris}#4}\par
  \vspace{1mm}
  \textbf{\lblFecha:} \rule{22mm}{0.4pt}\hspace{4mm}%
  \textbf{\lblLeido:} \casilla\ \lblLeidoSola\quad
  \casilla\ \lblLeidoConAyuda\quad
  \casilla\ \lblLeidoConDificultad\hspace{4mm}%
  \textbf{\lblNotas:} \rule{20mm}{0.4pt}\par
  \vspace{3mm}
}
\newenvironment{diapagina}[4]{%
  \def\diapagina@actual{#1}%
  \cabeceraDia{#1}{#2}{#3}{#4}%
}{%
  \vfill
  \begin{center}\footnotesize\color{colorGris}\lblDia~\diapagina@actual~/ \totaldias\end{center}
  \newpage
}
